% =========================================================
% mechlib/elements/supports.tex
% ATOMARE LAGER: FESTLAGER, LOSLAGER UND EINSPANNUNG
% =========================================================
% Diese Datei kapselt die Lagergeometrie. Fest- und Loslager teilen sich die
% Key-Familie /mech/support-common, damit Breite, Hoehe, Schraffur und Label
% konsistent parametrisiert werden koennen.
%
% Oeffentliche Pics:
%   mech/pinned-support Festlager / gelenkiges Lager
%   mech/roller-support Loslager mit sichtbarem Abstand zum Boden
%   mech/fixed-support  Einspannung, optional gedreht
%
% Aktuelle Fachwerk-Konvention:
%   - Die Dreiecksspitze liegt exakt im Mittelpunkt des Gelenks (0,0).
%   - Das Dreieck ist gegenueber der frueheren Version kompakt skaliert.
%   - Bodenlinie und Schraffur verwenden dieselben x-Grenzen wie die
%     Dreiecksgrundseite; beim Loslager wird die sichtbare Linienbreite explizit
%     beruecksichtigt.
%   - Das Gelenk selbst ist mit mechBodyColor gefuellt.
%
% Die pgfmathsetlengthmacro-Aufrufe sind absichtlich verwendet: LaTeX-Laengen
% duerfen nicht durch Textersatz wie ".5\\macro" skaliert werden, da dies bei
% expandierenden Dimensionsmakros zu sichtbarem Resttext (z. B. ".8mm") fuehren
% kann. Alle abgeleiteten Abmessungen werden deshalb numerisch berechnet.
% =========================================================
\pgfkeys{
  /mech/support-common/.is family,
  /mech/support-common,
  width/.initial=2.72mm,
  height/.initial=2.24mm,
  hatch angle/.initial=45,
  label/.initial=,
  label position/.initial=above,
}
\tikzset{
  % -----------------------------------------------------------------------
  % Festlager
  % -----------------------------------------------------------------------
  pics/mech/pinned-support/.style args={#1}{
    code={
      \pgfkeys{/mech/support-common/.cd,#1}
      \edef\mechSupportWidth{\pgfkeysvalueof{/mech/support-common/width}}
      \edef\mechSupportHeight{\pgfkeysvalueof{/mech/support-common/height}}
      \edef\mechSupportHatch{\pgfkeysvalueof{/mech/support-common/hatch angle}}
      \pgfmathsetlengthmacro{\mechSupportHalfWidth}{0.5*\mechSupportWidth}
      % Die Dreieckskontur ragt an der Grundseite durch die halbe Linienbreite
      % minimal ueber die geometrische Breite hinaus. Fuer identische sichtbare
      % Breite wird die Bodenlinie daher um 0.3pt je Seite erweitert.
      \pgfmathsetlengthmacro{\mechSupportGroundHalfWidth}{\mechSupportHalfWidth+0.3pt}
      \pgfmathsetlengthmacro{\mechSupportApexY}{0mm}
      \pgfmathsetlengthmacro{\mechSupportHatchDepth}{0.256*\mechSupportHeight}

      \coordinate (-connection) at (0,0);
      \coordinate (-center) at (0,0);

      \draw[mech/support]
        (-\mechSupportHalfWidth,-\mechSupportHeight)
        -- (0,\mechSupportApexY)
        -- (\mechSupportHalfWidth,-\mechSupportHeight)
        -- cycle;
      \draw[mech/body] (0,0) circle[radius=.6mm];
      \draw[black,line width=.6pt,line cap=butt]
        (-\mechSupportGroundHalfWidth,-\mechSupportHeight)
        -- (\mechSupportGroundHalfWidth,-\mechSupportHeight);
      \fill[
        pattern={Lines[angle=\mechSupportHatch,distance=1pt]},
        pattern color=black
      ]
        (-\mechSupportGroundHalfWidth,-\mechSupportHeight)
        rectangle
        (\mechSupportGroundHalfWidth,{-(\mechSupportHeight+\mechSupportHatchDepth)});

      \node[\pgfkeysvalueof{/mech/support-common/label position}]
        at (0,0) {\pgfkeysvalueof{/mech/support-common/label}};
    }
  },
  pics/mech/pinned-support/.default={},
  % -----------------------------------------------------------------------
  % Loslager
  % -----------------------------------------------------------------------
  pics/mech/roller-support/.style args={#1}{
    code={
      \pgfkeys{/mech/support-common/.cd,#1}
      \edef\mechSupportWidth{\pgfkeysvalueof{/mech/support-common/width}}
      \edef\mechSupportHeight{\pgfkeysvalueof{/mech/support-common/height}}
      \edef\mechSupportHatch{\pgfkeysvalueof{/mech/support-common/hatch angle}}
      \pgfmathsetlengthmacro{\mechSupportHalfWidth}{0.5*\mechSupportWidth}
      % The triangle outline extends by half its line width at the base corners.
      % Add 0.3pt so the roller ground has the same visible total width.
      \pgfmathsetlengthmacro{\mechSupportGroundHalfWidth}{\mechSupportHalfWidth+0.3pt}
      \pgfmathsetlengthmacro{\mechSupportApexY}{0mm}
      \pgfmathsetlengthmacro{\mechSupportRollerGap}{0.272*\mechSupportHeight}
      \pgfmathsetlengthmacro{\mechSupportHatchDepth}{0.256*\mechSupportHeight}
      \pgfmathsetlengthmacro{\mechSupportGroundY}{\mechSupportHeight+\mechSupportRollerGap}

      \coordinate (-connection) at (0,0);
      \coordinate (-center) at (0,0);

      \draw[mech/support]
        (-\mechSupportHalfWidth,-\mechSupportHeight)
        -- (0,\mechSupportApexY)
        -- (\mechSupportHalfWidth,-\mechSupportHeight)
        -- cycle;
      \draw[mech/body] (0,0) circle[radius=.6mm];
      \draw[black,line width=.6pt,line cap=butt]
        (-\mechSupportGroundHalfWidth,-\mechSupportGroundY)
        -- (\mechSupportGroundHalfWidth,-\mechSupportGroundY);
      \fill[
        pattern={Lines[angle=\mechSupportHatch,distance=1pt]},
        pattern color=black
      ]
        (-\mechSupportGroundHalfWidth,-\mechSupportGroundY)
        rectangle
        (\mechSupportGroundHalfWidth,{-(\mechSupportGroundY+\mechSupportHatchDepth)});

      \node[\pgfkeysvalueof{/mech/support-common/label position}]
        at (0,0) {\pgfkeysvalueof{/mech/support-common/label}};
    }
  },
  pics/mech/roller-support/.default={},
}
\pgfkeys{
  /mech/fixed-support/.is family,
  /mech/fixed-support,
  width/.initial=1cm,
  depth/.initial=1.5mm,
  hatch angle/.initial=45,
  angle/.initial=0,
  label/.initial=,
}
\tikzset{
  pics/mech/fixed-support/.style args={#1}{
    code={
      \pgfkeys{/mech/fixed-support/.cd,#1}
      \edef\mechFixedWidth{\pgfkeysvalueof{/mech/fixed-support/width}}
      \edef\mechFixedDepth{\pgfkeysvalueof{/mech/fixed-support/depth}}
      \edef\mechFixedHatch{\pgfkeysvalueof{/mech/fixed-support/hatch angle}}
      \edef\mechFixedAngle{\pgfkeysvalueof{/mech/fixed-support/angle}}

      \coordinate (-center) at (0,0);
      \coordinate (-connection-upper)  at (0,\mechConnectionOffset);
      \coordinate (-connection-middle) at (0,0);
      \coordinate (-connection-lower)  at (0,-\mechConnectionOffset);

      \begin{scope}[rotate=\mechFixedAngle]
        \draw[mech/body,fill=none] (-.5\mechFixedWidth,0)--(.5\mechFixedWidth,0);
        \fill[
          pattern={Lines[angle=\mechFixedHatch,distance={3pt/sqrt(2)}]},
          pattern color=black
        ]
          (-.5\mechFixedWidth,0)
          rectangle
          (.5\mechFixedWidth,\mechFixedDepth);
      \end{scope}

      \node[above] at (0,0) {\pgfkeysvalueof{/mech/fixed-support/label}};
    }
  },
  pics/mech/fixed-support/.default={},
}
